\section{Managerial, Policy, and Societal Implications}

Beyond its forecasting function, the cohort-based framework developed in this study operates as a diagnostic instrument that translates biological herd dynamics into decision-driven intelligence. By resolving each value chain into age and sex-structured cohorts advanced on a quarterly basis, its outputs speak simultaneously to herd and enterprise managers, to national and county policymakers, and to the broader societal objectives of food security, livelihood resilience, and environmental sustainability.

\subsection{Managerial Implications}
At the enterprise level, the framework identifies precisely where within a production system growth is being lost by isolating the cohorts and transition processes that constrain performance---e.g., elevated calf, lamb, and kid mortality, depressed reproductive rates, and sub-optimal off-take structures. This granularity matters because, as seen from comparable East African rangeland systems indicates, the sale and retention rates of breeding and mature females exert the greatest influence on long-run population trajectories \cite{tuffa2017, hary2004}. 

Managers may therefore use model outputs to assess herd structure, inform culling and destocking decisions, and breeding interventions. The quarterly time step also facilitates monitoring seasonal changes in herd dynamics and feed requirements, potentially supporting early intervention where projected feed deficits or changes in herd condition indicate emerging risks \cite{meshesha2019}. 

\subsection{Policy Implications}
For national and county governments the framework directly translates into the analytical needs of national and county planning providing the value chain specific projections necessary to prioritize investment and set performance targets. The ability to run coherent forward scenarios across nine value chains provides a basis for evidence-based investment sequencing and could be institutionalized within national statistical and livestock agencies, transforming sector planning from static census snapshots to continuous data-driven foresight \cite{bosire2022}.

\subsection{Societal Implications}
At the societal scale, the framework addresses the demand-driven transformation of animal-source-food consumption that will accompany Kenya’s population growth and urbanization. Reliable projections of milk, meat, egg, and honey supply against rising demand are central to food and nutrition security planning, allowing anticipatory action on emerging supply gaps rather than reactive crisis response. Since livestock underpins the livelihoods of several million pastoralists and agro-pastoralists and constitutes the dominant source of income across the arid and semi-arid lands, the model's capacity to anticipate feed shortfalls and productivity shocks has direct implications for poverty reduction, resilience, and the protection of vulnerable rural households \cite{behnke2011}.